% Préambule commun aux manuels de Pipit (français et anglais).
% Compilation : pdflatex manuel-fr.tex (deux fois), idem manuel-en.tex, depuis docs/src/.
\documentclass[11pt,a4paper]{article}
\usepackage[utf8]{inputenc}
\usepackage[T1]{fontenc}
\usepackage{lmodern}
\usepackage{microtype}
\usepackage[a4paper,margin=2cm,headheight=24pt]{geometry}
\usepackage{graphicx}
\usepackage[table]{xcolor}
\usepackage[most]{tcolorbox}
\usepackage{enumitem}
\usepackage{booktabs}
\usepackage{array}
\usepackage{tabularx}
\usepackage[nobottomtitles*]{titlesec}
\usepackage{needspace}
\usepackage{fancyhdr}
\usepackage{fontawesome5}
\usepackage{etoolbox}
\DeclareUnicodeCharacter{2192}{\ensuremath{\rightarrow}}
\DeclareUnicodeCharacter{2191}{\ensuremath{\uparrow}}
\DeclareUnicodeCharacter{2193}{\ensuremath{\downarrow}}
\DeclareUnicodeCharacter{25B2}{\ensuremath{\blacktriangle}}
\DeclareUnicodeCharacter{25BC}{\ensuremath{\blacktriangledown}}
\DeclareUnicodeCharacter{21BA}{\ensuremath{\circlearrowleft}}
\DeclareUnicodeCharacter{21BB}{\ensuremath{\circlearrowright}}
\usepackage{amssymb}

% Palette de l'application : vert olive, rouge brique pour les alertes
\definecolor{bordeaux}{HTML}{3A5420}
\definecolor{bordeauxl}{HTML}{EEF4E4}
\definecolor{alert}{HTML}{B5382B}
\definecolor{alertl}{HTML}{F8E6E3}
\definecolor{or}{HTML}{B8860B}
\definecolor{orl}{HTML}{FBF3DC}
\definecolor{sapin}{HTML}{2C7A55}
\definecolor{sapinl}{HTML}{E2F1E9}
\definecolor{sky}{HTML}{4E6E2A}
\definecolor{skyl}{HTML}{EEF4E4}
\definecolor{summit}{HTML}{6B7A8F}
\definecolor{night}{HTML}{2A2420}
\definecolor{ink}{HTML}{1F2430}
\AtBeginDocument{\color{ink}}

\titleformat{\section}
  {\Large\bfseries\sffamily\color{bordeaux}}
  {\makebox[1.6em][l]{\thesection}}{0.3em}{}
  [{\vspace{1pt}\color{bordeaux!35}\titlerule[1.3pt]}]
\titlespacing*{\section}{0pt}{14pt}{7pt}
\titleformat{\subsection}{\large\bfseries\sffamily\color{night}}{\thesubsection}{0.5em}{}
\pretocmd{\section}{\needspace{10\baselineskip}}{}{}
\pretocmd{\subsection}{\needspace{7\baselineskip}}{}{}
\widowpenalty=10000
\clubpenalty=10000
\raggedbottom
\setlength{\parindent}{0pt}
\setlength{\parskip}{6pt}
\setlist[enumerate]{itemsep=3pt,leftmargin=6mm}
\setlist[itemize]{itemsep=2pt,leftmargin=5mm}

\pagestyle{fancy}
\fancyhf{}
\renewcommand{\headrulewidth}{0.4pt}
\renewcommand{\headrule}{\hbox to\headwidth{\color{bordeaux!40}\leaders\hrule height \headrulewidth\hfill}}
\fancyfoot[C]{\sffamily\small\color{summit}\thepage}
\renewcommand{\sectionmark}[1]{\markboth{#1}{}}

\newcommand{\menu}[1]{\textsf{\textbf{\textcolor{bordeaux}{#1}}}}
\newcommand{\touche}[1]{\tcbox[on line,boxsep=0pt,left=2pt,right=2pt,top=1.5pt,bottom=1.5pt,
  colback=white,colframe=summit,boxrule=0.5pt,arc=1.5pt]{\footnotesize\textsf{#1}}}
\newcommand{\fichier}[1]{\texttt{\small #1}}

\newcommand{\boite}[5]{\newtcolorbox{#1}{enhanced,unbreakable,colback=#2,colframe=#3,
  coltitle=white,fonttitle=\bfseries\sffamily\small,boxrule=0.9pt,arc=2.5pt,
  left=2.5mm,right=2.5mm,top=1.2mm,bottom=1.2mm,title={#4~~#5}}}
\newcommand{\capture}[2][1]{\par\needspace{8\baselineskip}\begin{center}\setlength{\fboxsep}{0pt}\setlength{\fboxrule}{0.6pt}%
  \fcolorbox{summit!50}{white}{\includegraphics[width=\dimexpr#1\linewidth-1.3pt\relax]{#2}}\end{center}}

\usepackage[french]{babel}
\usepackage[colorlinks=true,linkcolor=bordeaux,urlcolor=bordeaux]{hyperref}
\hypersetup{pdftitle={Pipit — manuel d'utilisation},pdfauthor={},pdfcreator={Pipit}}
\boite{astucebox}{orl}{or!85!black}{\faLightbulb}{Bon à savoir}
\boite{securitebox}{sapinl}{sapin}{\faShield*}{Ce que Pipit ne fait jamais}
\boite{attentionbox}{alertl}{alert}{\faExclamationTriangle}{Attention}
\newcommand{\astuce}[1]{\par\needspace{5\baselineskip}\begin{astucebox}#1\end{astucebox}}
\newcommand{\attention}[1]{\par\needspace{5\baselineskip}\begin{attentionbox}#1\end{attentionbox}}
\fancyhead[L]{\sffamily\small\color{bordeaux}\textbf{PIPIT — MANUEL}}
\fancyhead[R]{\sffamily\small\color{summit}\leftmark}
\input{version}

\begin{document}

\begin{titlepage}
\centering
\vspace*{2.2cm}
\includegraphics[width=3.6cm]{../../logo_512.png}\par
\vspace{8mm}
{\sffamily\Huge\bfseries\color{bordeaux} PIPIT\par}
\vspace{3mm}
{\sffamily\LARGE\color{ink} Retirer une clé ou un disque USB sans risque\par}
\vspace{2mm}
{\sffamily\large\color{summit} Manuel d'utilisation \textbullet{} version \versionpipit\par}
\vspace{12mm}
\begin{tcolorbox}[enhanced,width=0.82\linewidth,colback=skyl,colframe=sky,arc=4pt,boxrule=1.2pt,left=5mm,right=5mm,top=3mm,bottom=3mm]
\sffamily\centering
Le pipit est un petit passereau discret des prés :\\
Pipit reste près de l'horloge et ne se montre que quand on l'appelle.\\[2mm]
\textbf{Clic sur l'icône} \quad → \quad \textbf{clic sur le périphérique} \quad → \quad \textbf{débrancher}
\end{tcolorbox}
\vfill
{\footnotesize\color{summit}\sffamily Windows 10 et 11 \textbullet{} logiciel libre (licence MIT)\par}
\end{titlepage}

\tableofcontents
\newpage

\section{À quoi sert Pipit}

Débrancher une clé USB ou un disque externe pendant que Windows écrit encore dessus peut
abîmer des fichiers, voire tout le disque. Windows a bien son icône « Retirer le
périphérique en toute sécurité », mais elle est cachée, ne dit pas \emph{quel} programme
bloque quand elle refuse, et n'a pas de raccourci clavier.

Pipit fait la même chose, en mieux rangé :

\begin{itemize}
\item une icône près de l'horloge : un clic montre la liste, un clic sur une ligne éjecte ;
\item un raccourci clavier, actif partout (\touche{Win}+\touche{Alt} par défaut) ;
\item quand Windows refuse, le nom du programme en cause, et un bouton pour le fermer ;
\item ni le disque de Windows, ni les disques internes ne sont jamais proposés.
\end{itemize}

Pipit est gratuit et le restera : pas de licence, pas de version « Pro », pas de
mise à jour payante.

\section{Installer}

\begin{enumerate}
\item Télécharger \fichier{Pipit-Setup.exe} sur la page
  \url{https://github.com/ARP273-ROSE/Pipit/releases/latest}.
\item Le lancer : l'installation se fait pour votre compte, sans fenêtre
  d'administrateur.
\item Au premier lancement, Pipit pose deux questions : envoyer ou non les rapports
  d'incident, et se lancer ou non avec Windows. Répondez \menu{Oui} à la seconde : c'est
  ce qui le rend toujours disponible quand vous branchez une clé.
\end{enumerate}

\astuce{Si l'icône n'apparaît pas près de l'horloge, cliquez sur la petite flèche
\textbf{\^{}} de la barre des tâches : Windows range là les icônes nouvelles. Faites
glisser celle de Pipit dans la barre pour l'avoir toujours sous la main.}

\section{Retirer un périphérique}\label{sec:retirer}

Cliquez sur l'icône de Pipit (le symbole « éjecter » vert) ou appuyez sur
\touche{Win}+\touche{Alt}. Le menu s'ouvre sous la souris :

\capture[0.62]{../capture_menu-fr.png}

Chaque ligne donne la lettre du lecteur, son nom et sa taille. \textbf{Un clic sur la
ligne éjecte le périphérique.} Quand c'est fait, une bulle « Vous pouvez débrancher »
apparaît : vous pouvez retirer le câble.

\begin{itemize}
\item \menu{Tout éjecter} retire tout ce qui est dans la liste, d'un coup.
\item \menu{Ouvrir dans l'Explorateur} ouvre le contenu d'un lecteur.
\item \menu{Réglages…} ouvre la fenêtre décrite section~\ref{sec:reglages}.
\end{itemize}

\section{Quand Windows refuse}\label{sec:refus}

Windows refuse d'éjecter tant qu'un programme garde un fichier ouvert sur le
périphérique : un document Word, une photo dans une visionneuse, une vidéo en pause…
Pipit refait d'abord un essai une seconde plus tard (un antivirus qui examine la clé la
lâche souvent tout seul), puis, s'il échoue encore, dit qui est en cause :

\capture[0.6]{../capture_refus-fr.png}

\begin{itemize}
\item \menu{Fermer …} demande poliment au programme de se fermer, exactement comme la
  croix de sa fenêtre : s'il a un document non enregistré, il vous le proposera. Pipit
  réessaie ensuite tout seul.
\item \menu{Réessayer} relance l'éjection, après que vous avez fermé vous-même ce qu'il
  fallait.
\end{itemize}

Cas particuliers :
\begin{itemize}
\item \textbf{l'Explorateur de fichiers} : une fenêtre est ouverte sur le lecteur, ou le
  volet d'aperçu affiche un de ses fichiers. Fermez cette fenêtre ; Pipit ne ferme jamais
  l'Explorateur, qui emporterait avec lui le bureau et la barre des tâches ;
\item \textbf{un antivirus ou l'indexation de Windows} : patientez quelques secondes et
  réessayez ;
\item \textbf{aucun programme nommé} : Windows ne l'a pas noté. Fermez les documents
  ouverts depuis le périphérique.
\end{itemize}

\begin{securitebox}
\begin{itemize}[leftmargin=4mm]
\item forcer une éjection : si Windows refuse, rien n'est retiré ;
\item fermer de force un programme (vos documents non enregistrés sont protégés) ;
\item fermer l'Explorateur, l'antivirus ou un programme de Windows ;
\item proposer le disque de Windows ;
\item demander des droits d'administrateur, installer un pilote ou un service.
\end{itemize}
\end{securitebox}

\section{Lecteurs de cartes à plusieurs fentes}

Un lecteur qui accepte à la fois une carte SD et une carte CF apparaît comme un seul
périphérique avec plusieurs lettres. Windows ne sait éjecter que le lecteur entier :
Pipit vous le rappelle avant de le faire.

\capture[0.6]{../capture_lecteur-fr.png}

Cochez \menu{Ne plus demander} si vous n'utilisez qu'une fente à la fois.

\section{Réglages}\label{sec:reglages}

\menu{Réglages…} dans le menu de l'icône ouvre cette fenêtre. La fermer ne quitte pas
Pipit : l'icône reste près de l'horloge. Pour quitter vraiment : \menu{Quitter Pipit}
dans le menu.

\capture{../capture_reglages-fr.png}

\subsection{Périphériques branchés}
\begin{itemize}
\item \textbf{Nom affiché} : double-cliquez pour donner un nom parlant
  (« Sauvegarde photo » plutôt que « SanDisk Extreme Pro USB Device »). Effacez le nom
  pour revenir à celui du matériel.
\item \textbf{Masquer} : le périphérique n'apparaît plus dans le menu (un disque que vous
  ne débranchez jamais, par exemple). Nom et masquage sont retenus par modèle de
  matériel : Pipit s'en souvient même quand le périphérique est débranché.
\end{itemize}

\subsection{Raccourcis clavier}
\begin{itemize}
\item \textbf{Ouvrir le menu} : \touche{Win}+\touche{Alt} par défaut (le même que dans
  USB Safely Remove), ou un autre au choix.
\item \textbf{Tout éjecter} : aucun par défaut, pour éviter une éjection par mégarde.
\end{itemize}
Si un autre programme utilise déjà le raccourci choisi, Pipit le signale sous la liste :
choisissez-en un autre.

\subsection{Options}
\begin{itemize}
\item \textbf{Prévenir quand on peut débrancher} : la bulle de confirmation.
\item \textbf{Demander avant d'éjecter un lecteur multicarte} : voir plus haut.
\item \textbf{Montrer aussi les disques internes branchés à chaud} : Windows présente
  souvent les disques SATA internes comme « amovibles ». Pipit les cache pour éviter
  d'éjecter un disque interne par erreur ; cochez seulement si vous avez un tiroir
  ou une baie extractible.
\item \textbf{Lancer Pipit au démarrage de Windows}.
\item \textbf{Langue} : français, anglais, ou celle de Windows.
\end{itemize}

\section{Mises à jour et rapports d'incident}

Pipit vérifie discrètement, au lancement, s'il existe une nouvelle version, et vous la
propose. Chaque mise à jour est \textbf{signée} : avant de l'installer, Pipit vérifie
qu'elle vient bien de son auteur et qu'elle n'a pas été modifiée en route.

Si Pipit plante ou se fige, et seulement si vous l'avez accepté, un rapport technique
part tout seul (version, système, trace de l'erreur ; ni vos fichiers, ni la liste de vos
périphériques). \menu{Signaler un problème…} permet d'écrire soi-même.

\section{Ligne de commande}

Depuis une invite de commandes, dans le dossier de Pipit :

\begin{center}
\begin{tabularx}{\linewidth}{>{\ttfamily\small}l X}
\toprule
\textrm{\textbf{Commande}} & \textbf{Effet} \\
\midrule
run.py --lister & liste ce qui peut être retiré, sans rien éjecter \\
run.py --ejecter E & éjecte le périphérique qui porte la lettre E: \\
run.py --discret & lance l'icône seule, sans ouvrir les réglages \\
\bottomrule
\end{tabularx}
\end{center}

\section{Questions fréquentes}

\textbf{Mon disque externe n'apparaît pas.} S'il est branché en eSATA ou dans une baie
SATA, cochez « Montrer aussi les disques internes branchés à chaud ». S'il est masqué,
décochez \menu{Masquer} dans les réglages.

\textbf{Une souris, un clavier, un téléphone apparaissent-ils ?} Non : Pipit ne liste que
les disques (clés, disques, cartes mémoire). Le reste se débranche sans précaution.

\textbf{Après l'éjection, le disque continue de tourner.} Selon le boîtier, l'alimentation
n'est pas coupée ; il est pourtant bien démonté. Vous pouvez débrancher.

\textbf{Je veux revoir un périphérique éjecté sans le débrancher.} Débranchez-le puis
rebranchez-le : Windows ne sait pas « réinsérer » un périphérique éjecté.

\end{document}
